
利用区块封装正文组件的意义在于, 您可以在正文的任何位置使用 \cs{MakeBlockIndex} 命令来快速生成某群组下全体同类型区块的索引. 区块索引同时配备用于交叉引用的锚点, 因此可以在其它位置\blclink{添加区块索引链接-U}以便于读者快速定位到此处.

得益于群组的树形层级结构, 针对某一群组生成的索引将自动包含其所有下级群组中的区块条目, 这些下级群组索引的锚点也会同步生成.

\UserCommand{生成区块索引}<\cs{MakeBlockIndex}>{
    \cs{MakeBlockIndex}
    \{\meta{范围路径}\}
    [\meta{区块类型列表}]
    <\meta{自定义参数}>
}[
    - \meta{范围路径} : 用于指定索引中区块所属的群组路径.
][
    - \meta{区块类型列表} : 用于指定索引中包含的常规区块类型, 各类型间通过 \texttt{,} 分隔, 若此参数为空则默认生成所有区块类型的索引.
    - \meta{自定义参数} : 在\seclink{导言区/样式设定}中使用该参数来实现各种复杂的排版逻辑.
]

范围路径是群组索引用于交叉引用的唯一标识符, 因此你\textcolor{red}{不能生成两个存在范围重叠的区块索引}, 即使它们可能使用了不同的区块类型列表, 若确需生成范围重叠的区块索引, 您可以使用下面的命令来创建不配备锚点的临时区块索引:

\UserCommand{生成临时区块索引}<\cs{MakeBlockIndex*}>{
    \cs{MakeBlockIndex*}
    \{\meta{区块索引范围路径}\}
    [\meta{区块类型列表}]
    <\meta{自定义参数}>
}

\textcolor{red}{请不要在区块或其它命令的参数中使用生成区块索引的命令}.

\newpage
